\section{Full-rank theta decay and the indefinite theta integral}
\label{sec:theta-decay-proof}

Keep the arithmetic base \(B\), its CM extension \(E/F\) and involution
\(\sigma\), its additive character \(\psi\), the embeddings
\(\tau_w:E\to\mathbb C\) above the infinite places \(w\) of \(F\), and
its distinguished split finite place \(v_1\) with a chosen embedding
\(\tau_1:E\to F_{v_1}\). Put \(D=[E:\mathbb Q]\).
Let \(X,W\) be the nondegenerate Hermitian and skew-Hermitian spaces
of dimensions \(n,d\) and signatures \((r,s),(r',s')\) in the pair
of Section~\ref{sec:theta-match}. Write \(U_X,U_W\) for their adelic
unitary groups, \(M\) for the specified Schr\"odinger model,
\(\omega\) for its Weil action, \(\Theta_\Phi(g,h)\) for its theta
sum, and \(\mathcal F_X\) for its rational fundamental domain with
the model's Haar measure \(\mu_X\). The pair hypothesis includes the
good base and model, the Hecke characters and their required restrictions,
and the Weil action on this model.

Let \(\mathcal S_\infty\) denote the Schwartz space on the infinite
part of \(\mathbb A_F^{nd}\), and \(\mathcal S_f\) the compactly
supported, locally constant functions on the finite part. In the
chosen split chart, \(M.\mathrm{SoLoc}\) consists of local
Schwartz--Bruhat functions whose image vanishes on matrices of row
rank less than \(n\). The image is taken under a choice of admissible local
intertwiner in the definition of the chart. In the global space
\(M.\mathrm{So}\) take the span of pure tensors with this condition at
\(v_1\), arbitrary Schwartz factors elsewhere and the specified
vacuum at almost every finite place. Set
\[
 \mathcal S_f^\circ=
 \{\,\phi\in\mathcal S_f:
          f\otimes\phi\in M.\mathrm{So}\quad
               \text{for every } f\in\mathcal S_\infty\,\}.
\]
Each member of \(M.\mathrm{So}\) is a finite sum of such tensors.
Let \(H\) be the matrix height of the model and
\(H_Q(g)=\inf_{\gamma\in U_X(F)}H(\gamma g)\).

For each \(N_0\in\mathbb N\) there is \(k\in\mathbb N\) such that
for each \(\phi\in\mathcal S_f^\circ\) there is a continuous seminorm
\(p_\phi\) on \(\mathcal S_\infty\) satisfying
\begin{equation}
 |\Theta_{f\otimes\phi}(g,h)|H_Q(g)^{N_0}
       \le p_\phi(f) H(h)^k
 \qquad(f\in\mathcal S_\infty,\ g\in U_X,\ h\in U_W).
 \label{eq:theta-full-decay}
\end{equation}
The exponent \(k\) is uniform in \(\phi\). For unrestricted finite
Schwartz data, the hypothesis
\begin{equation}
 d+2\min(r,s)<2n\quad\text{or}\quad\min(r,s)=0
 \label{eq:theta-weil-range}
\end{equation}
gives a possibly different uniform \(k\) and, for each \(\phi\),
a continuous seminorm \(p_\phi\) and a nonnegative measurable
\(G_\phi\in L^1(\mathcal F_X,\mu_X)\) such that
\[
 |\Theta_{f\otimes\phi}(g,h)|
       \le p_\phi(f)H(h)^kG_\phi(g).
\]
We derive these statements with the same local data and measure.

\subsection{The square of the kernel}

For \(z\in\operatorname{Mat}_{n,d}(\mathbb A_F\otimes_F E)\)
write \(|z|_\infty=\sum_{w,i,j}|\pi_{\tau_w}(z_{ij})|\). Denote its
finite projection by \(z_f\). For compact \(C\) in this matrix space put
\[
 b_{C,N}(z)=\mathbf 1_{\{z_f\in C_f\}}\,(1+|z|_\infty)^{-N},
 \qquad
 L_{C,N,Y}(g,h)=\sum_{x\in Y} b_{C,N}(g^{-1}x(h^{-1})^t).
\]
Here \(C_f\) is the projection of \(C\), rational matrices are
embedded diagonally, and \(Y\) is either all rational matrices or
\(Y_n=\{x:\operatorname{rank}_E x=n\}\).
For each \(N\) and \(\phi\in\mathcal S_f\), the coefficient estimate produces
a compact \(C\) and a continuous seminorm \(p\) with
\begin{equation}
 |\Theta_{f\otimes\phi}(g,h)|^2\le p(f)^2 L_{C,N,Y}(g,h)
 \label{eq:theta-lattice-majorant}
\end{equation}
for \(Y\) all matrices. For \(\phi\in\mathcal S_f^\circ\) it produces
the same conclusion with \(Y=Y_n\).

Here are the local cancellation and the role of the square in the
second assertion. Write \(P\Phi(x)=\Phi(-x)\), and let
\[
 c_\Phi(w)=\int(\operatorname{Schr}_\psi(w)\Phi)(u)
                    \overline{\Phi(-u)}\,du .
\]
If \(\operatorname{Schr}_\psi(v)\Phi=\psi(a)\Phi\), translation
invariance of additive Haar measure gives
\[
 c_\Phi(w)=\psi(v_1\cdot w_2-w_1\cdot v_2+2a)c_\Phi(w).
\]
The subscripts on \(v,w\) in this identity denote their two
polarization coordinates. Indeed translation of the integration
variable by \(v_1\) contributes the commutator of the two Schr\"odinger
operators, and reflection followed by conjugation contributes the
second copy of \(a\). The parity operator on the theta sum gives
\(\Theta_{P\Phi}(g,h)=\Theta_\Phi(g,h)\), by covariance and the
bijection \(x\mapsto-x\) on the rational summation set. The product
bound for this pair is thus a bound for an absolute square.

In the split chart take a finite subdivision of the support of each
local full-rank factor into compact open pieces. The source's local
separation lemma chooses the pieces so that for every rank-deficient
matrix in the complementary split coordinate there is a Schr\"odinger
operator acting on the piece by a scalar for which the displayed
multiplier differs from \(1\). The complementary coordinate is
multiplied on the left and right by the invertible matrices of the
two forms and by a sign. Its rank is preserved. The coefficient
therefore vanishes there. For a rational matrix of rank less than
\(n\), extending scalars to this component and multiplying by
\(g^{-1},(h^{-1})^t\) still gives rank less than \(n\).
After this cancellation, the rational coefficient-sum formula,
Weil covariance and the Schwartz coefficient bound give
\eqref{eq:theta-lattice-majorant} on each piece. If \(p_0\) is
the seminorm for the two coefficients, use
\(p_0(f)+p_0(Pf)\) to bound their geometric mean. Formula
\eqref{eq:theta-lattice-majorant} is stable under finite sums. Take square roots, replace the finitely many compact sets by their
union, then add the seminorms with the absolute values of the scalar
coefficients. This explains the use of the full-rank locus after
taking the square. The remaining analytic inputs at this point are the coefficient-sum
formula, covariance and the Fourier seminorm estimate.

For completeness the separating subgroup has an explicit construction.
In \(k=F_{v_1}\), with valuation ring \(\mathcal O\), choose \(b\)
with \(\psi_{v_1}(b)\ne1\). Write
\(\langle Y,U\rangle=\operatorname{tr}(YU^t)\).
Let \(T\) be the compact nonzero set of the local factor in the split
chart. Compactness here follows from local constancy and compact
support. The full-rank locus is open, since its complement is the
common zero locus of all \(n\)-row minors. Since \(2\) is nonzero
in \(k\), \(2T\) is contained in the full-rank locus. Compactness
gives a neighborhood \(V_0\) of zero such that \(2T+V_0\) remains
in that locus. Take \(a\ne0\) with
\(a\operatorname{Mat}_{n,d}(\mathcal O)\subset V_0\) and put
\[
 K=(b/a)\operatorname{Mat}_{n,d}(\mathcal O),\qquad
 \Lambda=\{Y:\psi_{v_1}(\langle Y,U\rangle)=1\ (U\in K)\}.
\]
This subgroup is open by local constancy of the character, continuity
of the pairing and compactness of \(K\). It lies in
\(a\operatorname{Mat}_{n,d}(\mathcal O)\). Indeed if some
\(Y_{ij}/a\notin\mathcal O\), its inverse belongs to \(\mathcal O\);
taking \(U\in K\) with the single entry
\((b/a)(Y_{ij}/a)^{-1}\) gives \(\langle Y,U\rangle=b\).

The image of \(T\) in the discrete quotient by \(\Lambda\) is
finite. Restrict the function to each coset meeting \(T\), choose
\(Y_0\) in its intersection with \(T\), and let \(Z\) be rank deficient.
Since \(2Y_0+\Lambda\) is full rank, \(Z-2Y_0\notin\Lambda\).
Some \(U\in K\) thus has
\(\psi_{v_1}(\langle Z-2Y_0,U\rangle)\ne1\). On this whole
piece \(\psi_{v_1}(\langle Y,U\rangle)
=\psi_{v_1}(\langle Y_0,U\rangle)\).
The split polarization realizes this multiplication operator as the
Schr\"odinger action of the local vector with split coordinates
\((-U,0)\). Its eigenvalue is the common value just computed, and
the symplectic pairing with the complementary coordinate \(Z\)
gives \(\langle Z,U\rangle\) in the sign convention of the model.
The coefficient identity gives multiplier
\(\psi_{v_1}(\langle 2Y_0-Z,U\rangle)\ne1\), proving the vanishing.
This also covers residue characteristic two, since the argument uses
\(2\ne0\) in the characteristic-zero field and never divides by
\(2\) in the residue field.


\subsection{The two logarithmic chambers}

Reduction theory for unitary groups supplies rational Witt frames
of lengths \(q\le\min(r,s)\) and \(y\le\min(r',s')\), tori
\(a_X,a_W\), and compact sets \(\Omega_X,\Omega_W\). There are
rational \(\gamma,\delta\) such that
\(\gamma g=a_X(u)\kappa\), \(\delta h=a_W(v)\kappa'\) with
\(\kappa,\kappa'\) in these compact sets and
\[
 u_1\ge\cdots\ge u_q\ge0,\qquad v_1\ge\cdots\ge v_y\ge0.
\]
Here \(u,v\) henceforth denote real tuples. Put \(U=\sum u_i\),
\(V=\sum v_j\), and
\[
 B(u,v)=\max_{0\le i\le q}
   \left(\sum_{a=1}^i u_a-\sum_{b=1}^{\min(i,y)}v_b\right).
\]
The choices \(i=0,q\) give \(B\ge0\) and \(U\le V+B\). The frame
weights \(\lambda_i\) are the \(+u_i,-u_i\) and \(n-2q\) zeros.
The weights \(\lambda'_j\) are the \(+v_j,-v_j\) and \(d-2y\) zeros.
The positive weight sum is
\begin{align*}
 P(u,v)&=\sum_{i,j}\max(0,\lambda_i+\lambda'_j)\\
 &=2\sum_{i=1}^{q}\sum_{j=1}^{y}\max(u_i,v_j)
                  +(d-2y)U+(n-2q)V\\
 &\le dU+nV .
\end{align*}
Indeed each block of four signed pair sums contributes
\(u_i+v_j+|u_i-v_j|\). The frame height bounds and boundedness on
compact sets give constants \(A,A',k_0,k_1\) independent of \(g,h\) with
\[
 H_Q(g)\le A'e^{k_0U},\qquad e^V\le A H(h)^{k_1}.
\]
For the second inequality apply the inverse torus height bound
to the rational translate \(\delta^{-1}a_W(v)=h(\kappa')^{-1}\)
and use submultiplicativity and compactness. For the first use the
rational representative \(\gamma g\). The lattice sums are invariant
under rational left translation in either group. Right translation
over the two compact sets can multiply a bound by a constant and
enlarge \(C\). These constants are independent of \(u,v,f\).

\subsection{Counting the rational matrices}

Take rational Witt bases \(b,b'\), write \(\xi=b^{-1}x(b'^{-1})^t\),
and put \(z=a_X(u)^{-1}x(a_W(v)^{-1})^t\).
There is an integer \(m\ge1\) and a constant \(C_0\ge1\) such that
\[
 z_f\in C_f\ \Longrightarrow\ m\xi_{ij}\in\mathcal O_E,\qquad
 \sum_{i,j,w} e^{-\lambda_i-\lambda'_j}|\tau_w(\xi_{ij})|
                    \le C_0|z|_\infty .
\]
The finite tori have identity finite components, so the finite
coordinates of \(\xi\) lie in the image of a fixed compact set
under a fixed change of basis. A compact finite adelic set has a
common denominator, by a finite subcover of the sets of bounded
denominators. At infinity the two fixed basis changes have bounded
operator norms and the tori act diagonally, giving the second bound.

If \(\Lambda\) is a discrete subgroup of a real normed vector
space of dimension \(D\), then for \(L>D\) an integer,
\begin{equation}
 \sum_{\ell\in\Lambda}(1+e^{-a}\|\ell\|)^{-L}
       \le K_\Lambda e^{D\max(0,a)} .
 \label{eq:theta-onevar}
\end{equation}
Choose \(\epsilon>0\) such that the open balls of radius \(\epsilon\)
about the lattice points are disjoint. Put \(c=e^{-\max(0,a)}\).
For \(y\) in the ball about \(\ell\),
\(1+c\|y\|\le(1+\epsilon)(1+c\|\ell\|)\).
The summand in \eqref{eq:theta-onevar} is bounded by
\((1+c\|\ell\|)^{-L}\). Integrate the latter bound on each ball,
sum and divide by the volume of a ball. The result is at most
\[
 \frac{(1+\epsilon)^L}{\operatorname{vol}(B_\epsilon)}
 c^{-D}\int (1+\|y\|)^{-L}\,dy .
\]
The integral converges since the volumes of shells of radii \(2^j\)
are \(O(2^{jD})\). This proves \eqref{eq:theta-onevar}.

Choose \(N_b=(D+1)nd\). Set
\(A_{ij}=e^{-\lambda_i-\lambda'_j}\sum_w|\tau_w(\xi_{ij})|\).
The coordinate estimate gives \(1+A_{ij}\le C_0(1+|z|_\infty)\).
Multiplying these inequalities and inverting gives
\[
 (1+|z|_\infty)^{-N_b}\le
       C_0^{N_b}\prod_{i,j}(1+A_{ij})^{-(D+1)} .
\]
The rational change of basis is injective. Enlarge the summation
to all matrices with independent entries in \(m^{-1}\mathcal O_E\)
and use Tonelli's theorem to factor the nonnegative sum. In the
Minkowski embedding the norm of an element is at most
\(\sum_w|\tau_w(\xi_{ij})|\) for the supremum norm on the archimedean
coordinates, since the two embeddings over any real place have the same
absolute value. Apply \eqref{eq:theta-onevar} in each coordinate,
with \(a=\lambda_i+\lambda'_j\). The result is
\begin{equation}
 L_{C,N_b,\mathrm{all}}(a_X(u),a_W(v))\le K_C e^{DP(u,v)}.
 \label{eq:theta-all-torus}
\end{equation}

Here is the gain from full row rank. For \(x\in Y_n\) with
\(z_f\in C_f\), choose \(i\) attaining the maximum \(B(u,v)\).
The \(i\) rows of \(\xi\) with weights \(-u_1,\ldots,-u_i\)
are linearly independent, so some \(i\) distinct columns
\(c_1,\ldots,c_i\) give a square minor \(T_0\) with nonzero
determinant. The determinant \(m^i\det T_0\) is a nonzero algebraic
integer. Its absolute field norm is an integer at least \(1\),
so at an embedding \(\tau_w\) we have
\(m^i|\tau_w(\det T_0)|\ge1\). Use that embedding in the matrix
\[
 T_{ab}=e^{u_a-\lambda'_{c_b}}\,\tau_w((T_0)_{ab}).
\]
The sum of the chosen \(i\) distinct column weights is at most
\(\sum_{b\le\min(i,y)}v_b\), since the only positive weights
available are \(v_1,\ldots,v_y\) in decreasing order. Consequently
\[
 e^{B(u,v)}
 \le m^i|\det T|
 \le i!\big[m C_0(1+|z|_\infty)\big]^i
 \le \big[n!m C_0(1+|z|_\infty)\big]^{n+1}.
\]
The middle inequality follows by bounding the \(i!\) terms of the
determinant individually using the coordinate estimate above.
The empty minor at \(i=0\) has determinant \(1\). With
\(c=(n!m C_0)^{-1}\), we have proved
\(c\,e^{B/(n+1)}\le1+|z|_\infty\).
For any \(M_1\in\mathbb N\), raise the last inequality to
\(M_1(n+1)\), invert it, and use \eqref{eq:theta-all-torus} on
the remaining lattice sum. Since all summands are nonnegative, its
restriction to \(Y_n\) can only decrease it. Thus
\begin{equation}
 L_{C,N_b+M_1(n+1),Y_n}(a_X(u),a_W(v))
       \le K_{C,M_1}\exp\{DP(u,v)-M_1 B(u,v)\}.
 \label{eq:theta-full-torus}
\end{equation}
The exponent needed on the left is chosen before the compact set
\(C\), as required when applying the Schwartz estimate.

\subsection{From the chamber to an integrable majorant}

Take \(M_1=Dd+2k_0N_0\). Applying
\eqref{eq:theta-lattice-majorant} and \eqref{eq:theta-full-torus}
and taking square roots, the exponent after multiplication by
\(H_Q(g)^{N_0}\) is at most
\begin{align*}
 \tfrac12(DP-M_1B)+k_0N_0U
 &\le \tfrac12\{M_1(U-B)+DnV\}\\
 &\le \tfrac12(M_1+Dn)V
 \le (M_1+Dn)V .
\end{align*}
The height bound on \(e^V\) proves \eqref{eq:theta-full-decay}
with \(k=k_1(M_1+Dn)\). Absorb all fixed compact-set and height
constants into the continuous seminorm. In particular \(N_0=0\)
permits \(G_\phi=1\). The rational quotient of this unitary group
has finite Haar volume. For fixed \(h\) this proves absolute
integrability and gives the bound
\[
 \int_{\mathcal F_X}|\Theta_{f\otimes\phi}(g,h)|\,d\mu_X(g)
       \le\mu_X(\mathcal F_X)p_\phi(f)H(h)^k .
\]

For unrestricted finite data use
\(G_\phi(g)=L_{C',N_b,\mathrm{all}}(g,1)^{1/2}\).
To see why this function suffices, reduce the second group into
\(a_W(v)\Omega_W\). Each diagonal archimedean stretch in
\(a_W(v)\), or its inverse, has operator norm at most \(e^V\)
in Witt coordinates. It is the identity at the finite places.
The fixed changes of basis therefore give
\[
 L_{C,N_b,\mathrm{all}}(g,h)\le K e^{N_bV}
                L_{C',N_b,\mathrm{all}}(g,1).
\]
Use \(e^V\le A H(h)^{k_1}\) and take a square root in
\eqref{eq:theta-lattice-majorant}. This gives the asserted majorant
with \(k=k_1N_b\), as in the source, since \(H(h)\ge1\).
The indicator of the finite projection of a compact set is Borel,
and the summands are measurable; \(G_\phi\) is measurable.

It remains to integrate over the first group. For a nonnegative
measurable rationally left-invariant function majorized on
\(a_X(u)\Omega_X\) by \(b(u)\), the Siegel integration inequality gives
\[
 \int_{\mathcal F_X}G(g)\,d\mu_X(g)
 \le K\int_{u_1\ge\cdots\ge u_q\ge0} b(u)e^{-D\rho_n(u)}\,du,
 \qquad
 \rho_n(u)=\sum_{i=1}^q(n-2i+1)u_i .
\]
The compact set and constant here may depend on \(\mu_X,\mathcal F_X\).
Equation \eqref{eq:theta-all-torus} at \(v=0\) gives
\(b(u)=K'\exp(DdU/2)\) for \(G_\phi\). Use the differences
\(z_j=u_j-u_{j+1}\) with \(u_{q+1}=0\). This takes the chamber to
the positive orthant and has Jacobian \(1\). The cumulative coefficient
of \(z_j\) in \(D(dU/2-\rho_n(u))\) is
\[
 D\sum_{i=1}^j(d/2-n+2i-1)=Dj(d/2-n+j).
\]
Equation \eqref{eq:theta-weil-range} implies \(q=0\) or
\(d+2q<2n\), because \(q\le\min(r,s)\). Thus every coefficient
for \(j\ge1\) is strictly negative, and with product Lebesgue measure
the chamber integral, omitting its fixed prefactor, equals
\begin{equation}
 \prod_{j=1}^q \frac{1}{Dj(n-d/2-j)}.
 \label{eq:theta-integral-product}
\end{equation}
In rank zero the empty product is \(1\). This completes the
integrability calculation for unrestricted finite data. Decomposing a
global Schwartz function into finitely many pure tensors proves the
same convergence for it.

In the triple of dimensions \(a,a',a+a'\), the large first space
has signature \((a-p,p)\), the smaller one has \((q',a'-q')\),
and their common second space has signature
\((a+a'-(p+q'),p+q')\). The triple hypotheses give
\(a+a'+2\min(a-p,p)<2a\). Apply the unrestricted bound to the
large first space and the full-rank bound to the small first space
with its prescribed split test data. The pointwise constants depend
on the arithmetic model, on the chosen compact sets and on the
Schwartz data. The estimates above supply the convergence wherever
the family supplies these hypotheses; they give no uniform bound
on those constants as the arithmetic model varies.

\subsection{Decay in the variable left after integration}

For the later cuspidal member of the theta span, decay is also
needed in the remaining variable \(h\in U_W\). The source applies
the cuspidal part of unitary reduction theory under the following
hypotheses on a function \(A\) on \(U_W\). The function is left
invariant under rational points and right invariant under an open
finite level. It is archimedean smooth and lies in a family closed
under differentiation along real one-parameter subgroups. All
members of this family have polynomial growth with a common
exponent, with individual constants allowed. For each nonzero
rational totally isotropic subspace the constant term over the
associated unipotent radical vanishes. In coordinates, one can
test these subspaces on matrices \(C\) of full column rank with
\(C^*JC=0\), where \(J\) is the Hermitian matrix representing the
group. The column span and its orthogonal complement specify the
unipotent radical. For the skew-Hermitian group the source first
identifies it with a Hermitian unitary group by scaling the
conjugate matrix of its form by a nonzero scalar of the appropriate
sign under \(\sigma\).

Under these hypotheses the cuspidal estimate in the imported
reduction theorem gives a constant \(C_N\) on a Siegel set for each
\(N\). For an arbitrary \(h\), choose a rational translate
\(\delta h\) in that set. Left invariance and the definition of the
quotient height give
\[
 |A(h)|H_Q(h)^N\le |A(\delta h)|H(\delta h)^N\le C_N .
\]
In \texttt{SigRapid.sigRapid\_of}, invariance, smoothness and
uniform moderate growth are supplied by \texttt{sig\_nice} for
members of the specified theta span. Vanishing of the constant
terms is a separate hypothesis of this implication.
\texttt{CuspRapidOf.rapidG} and \texttt{CuspRapidOf.cuspRapid\_of}
make the passage from the Siegel set to the rational quotient.
This use of the cuspidal reduction theorem and the explicit
full-rank estimate \eqref{eq:theta-full-decay} concern different
variables in the iterated integrals.

\paragraph{Source and inputs.}
In the pinned source tree, under
\srcid{Development/WeilBound/ThetaLifts/FrameEstimates},
\srcid{SquareIntegrableWeilRange} defines the lattices and chambers.
The local separation and the square estimate are
\texttt{HermitianFormDetNeZero.soCut\_of},
\texttt{ThetaSumNeg.qb\_piece} and \texttt{ThetaSumNeg.sqSo\_of'}.
The one-coordinate packing, unrestricted sum, minor and full-rank
sum are \texttt{OneVar.oneVar\_of}, \texttt{LatTor.latTor\_of},
\texttt{ExistsColsDetNeZero.minor\_of} and
\texttt{LatFull.latFull\_of}. The two global estimates are
\texttt{SoSemi.soSemiU\_of} and \texttt{DomWeil.domWeilU\_of},
assembled in \srcid{WeilConv}. These names combine the identifying
file with the declaration. The imported reduction theorem
\srcid{GlobalTheory.Published.SiegelSetUnitary} supplies the cover,
height comparisons and integration inequality.
\srcid{FundamentalDomainFiniteVolume} supplies finiteness of the
model's quotient measure. The local coset argument appears in \texttt{RankIffRight.exists\_cut}.
The Heisenberg product formula and its Schwartz coefficient estimate, and the construction
of each arithmetic pair with these hypotheses, remain upstream
inputs to this chapter.
